% =========================
% TIPOGRAFIA
% =========================

\usepackage{microtype}
\usepackage{lettrine}

% =========================
% CORES INFOSOC (Cubismo laranja creme azul)
% =========================

\usepackage{xcolor}
\usepackage{pagecolor}
\usepackage{afterpage}
\usepackage{etoolbox}

% Definição exata da paleta de apoio
\definecolor{laranjaqueimado}{HTML}{A33B14}
\definecolor{laranjavibrante}{HTML}{FF3B00}
\definecolor{laranja}{HTML}{F26D21}
\definecolor{pessego}{HTML}{FFC299}
\definecolor{creme}{HTML}{FDFBF7}
\definecolor{bege}{HTML}{D9CBB8}
\definecolor{azul}{HTML}{2A75D3}
\definecolor{ciano}{HTML}{00C4D9}
\definecolor{azulcobalto}{HTML}{0047AB}
\definecolor{textoescuro}{HTML}{1A1A1A}

% =========================
% KOMA - FONTES E CORES ESTRUTURAIS
% =========================

%Todos os títulos, incluindo o sumário
\setkomafont{disposition}{\sffamily\color{azulcobalto}}

% Título da primeira página do livro
\setkomafont{title}{\sffamily\color{laranjavibrante}}
\setkomafont{author}{\sffamily}
\setkomafont{date}{\sffamily}



%Numeração da parte
\setkomafont{partnumber}{\sffamily\bfseries\color{laranjaqueimado}\Huge}
\setkomafont{part}{\sffamily\bfseries\color{laranjavibrante}\fontsize{60}{76}\selectfont}

% Capítulo é definido mais abaixos

% Seções
\setkomafont{section}{\sffamily\bfseries\color{laranjavibrante}\Large}
\setkomafont{subsection}{\sffamily\color{azul}}
\setkomafont{subsubsection}{\sffmaily\color{azulcobalto}}

% Seções no TOC
\usepackage{tocbasic}

% Comando auxiliar que recebe o texto da entrada como #1

\newcommand{\formatacaoparte}[1]{%
  \setlength{\fboxsep}{6pt}%
  \colorbox{pessego}{%
    \makebox[\dimexpr\linewidth-2\fboxsep+28pt\relax][l]{%
      \color{laranjaqueimado}\sffamily\bfseries\large #1%
    }%
  }%
}

\DeclareTOCStyleEntry[
  entryformat=\formatacaoparte,
  pagenumberformat={\color{laranjaqueimado}\sffamily\bfseries},
  linefill={\leavevmode\color{laranjaqueimado}\TOCLineLeaderFill},
  beforeskip=1.5em
]{tocline}{part}


\DeclareTOCStyleEntry[
  entryformat={\color{laranjaqueimado}\sffamily},
  pagenumberformat={\color{laranjaqueimado}\sffamily\bfseries},
  linefill=\TOCLineLeaderFill
]{tocline}{chapter}

\DeclareTOCStyleEntry[
  entryformat={\color{laranjavibrante}\sffamily},
  pagenumberformat=\sffamily,
  linefill=\TOCLineLeaderFill
]{tocline}{section}

\DeclareTOCStyleEntry[
  entryformat={\color{azul}\sffamily},
  pagenumberformat=\sffamily
]{tocline}{subsection}


% --
% =========================
% CAPÍTULOS
% =========================

\RedeclareSectionCommand[
  beforeskip=2.0cm,
  afterskip=1.0cm
]{chapter}

% Remove o ponto após o número (padrão KOMA é "1.")
\renewcommand*{\chapterformat}{\thechapter}

% #2 = número formatado | #3 = título
% \ifblank: se #2 vazio (cap. não numerado), não exibe a caixa
\renewcommand*{\chapterlinesformat}[3]{%
  \noindent
  \ifblank{#2}{}{%
    {%
      \setlength{\fboxsep}{0.45em}%
      \colorbox{bege}{%
        \hspace{0.3em}%
        {\color{white}\bfseries\huge%
          \rule[-0.25em]{0pt}{1.2em}#2}%
        \hspace{0.3em}%
      }%
    }%
    \par\vspace{0.7em}%
  }%
  {\sffamily\bfseries\Huge\color{laranjaqueimado}#3}%
  \par
}
% --

% =========================
% PARTES
% =========================

\renewcommand*{\partformat}{\partname~\thepart}
\renewcommand*{\partpagestyle}{empty}

% \pretocmd: ativa a cor de fundo ANTES da página de parte ser composta.
% Utiliza o pêssego da nova paleta para um fundo de transição suave.
\pretocmd{\part}{%
  \cleardoublepage
  \pagecolor{pessego}%
}{}{}

% \partheadendvskip: chamado DENTRO da página de parte, ao final.
% \afterpage{\nopagecolor} é acionado após a página ser finalizada,
% restaurando o fundo branco puro para as páginas seguintes.
\renewcommand*{\partheadstartvskip}{\null\vfill}
\renewcommand*{\partheadendvskip}{%
  \vfill\null
  \afterpage{\nopagecolor}%
}

% =========================
% CABEÇALHO
% =========================

\usepackage{scrlayer-scrpage}
\clearpairofpagestyles

\automark[section]{chapter}

\ihead{\headmark}
\ohead{\pagemark}

% Elementos de navegação em azul cobalto garantem sobriedade e leitura eficiente
\setkomafont{pagehead}{\sffamily\small\color{laranja}}
\setkomafont{pagenumber}{\sffamily\small\color{laranja}}

\KOMAoptions{headsepline=true}
% Colorir a linha do cabeçalho
\addtokomafont{headsepline}{\color{bege}}

% =========================
% LEGENDAS (Figuras e Tabelas)
% =========================

% Altera a cor e a fonte do texto da legenda
\setkomafont{caption}{\sffamily\small\color{azulcobalto}}

% Altera a cor e a fonte do rótulo ("Figura 1:", "Tabela 1:")
% O uso do \sffamily e \bfseries cria contraste com o texto descritivo
\setkomafont{captionlabel}{\sffamily\small\color{azulcobalto}}


% =========================
% NOTAS DE RODAPÉ
% =========================

% Altera a cor do texto da nota de rodapé
\setkomafont{footnote}{\color{azulcobalto}}

% Altera a cor do número/marcação da nota no rodapé da página
\setkomafont{footnotelabel}{\color{azulcobalto}}

% Opcional: Altera a cor do número sobrescrito no meio do texto principal
% \setkomafont{footnotereference}{\color{azulcobalto}}


% =========================
% NOTAS MARGINAIS (Integração Quarto)
% =========================

% 1. Assume o controle sobre as notas geradas por marcações do Quarto (.aside)
\usepackage{marginnote}
\renewcommand*{\marginfont}{\sffamily\small\color{azulcobalto}}

% 2. Mantém a intercepção do comando KOMA-Script para marcações manuais no texto
\let\oldmarginpar\marginpar
\renewcommand{\marginpar}[1]{\oldmarginpar{\sffamily\small\color{azulcobalto}#1}}

\let\oldmarginline\marginline
\renewcommand{\marginline}[1]{\oldmarginline{\sffamily\small\color{azulcobalto}#1}}